\subsection*{\texttt{\_\_device\_\_} functions}

When an algorithm needs a helper function that operates at the lane or thread level --- a custom reduction, a sorting network, a special math function --- writing it as a device function is often the natural choice. Croqtile does not need to manage these; they are passed through to the target compiler.

A \texttt{\_\_co\_\_} function can call a \texttt{\_\_device\_\_} function using the \texttt{call} keyword:

\begin{lstlisting}[language=croqtile,caption={device functions}]
#define N_TOKENS 1024
#define N_EXPERTS 8
#define N_THREAD 128
template <int K>
__device__ void warp_topk(f32* vals, s32* idxs);
template <typename T>
__device__ __forceinline__ T SHFL_XOR(T var, int lane_mask, int width) {
  return __shfl_xor_sync(uint32_t(-1), var, lane_mask, width);
}
__co__ void moe_topk(f32 [N_TOKENS, N_EXPERTS] scores,
                     s32 [N_TOKENS, K]& topk_ids,
                     f32 [N_TOKENS, K]& topk_scores,
                     int N_BLOCK) {
  parallel n by N_BLOCK : block {
    foreach m in [ |scores.span| / N_THREAD / N_BLOCK ] {
      shared_scores = dma.copy scores.chunkat(n#m, _) => shared;
      parallel gid by [N_THREAD / 32] : group
        parallel tid by 32 : thread {
          flat_scores = shared_scores.span_as([|shared_scores.span|]);
          score = flat_scores.at(gid # tid);
          local s32 [K] frag_idx{-1};
          local f32 [K] frag_val{-1.0f};
          frag_idx.at(0) = tid;
          frag_val.at(0) = score;
          call warp_topk<8>(frag_val, frag_idx);
          inthreads (tid == 0)
            foreach k in K {
              topk_ids.at(n#m, gid#k) = frag_idx.at(k);
              topk_scores.at(n#m, gid#k) = frag_val.at(k);
            }
        }
    }
  }
}
\end{lstlisting}

\textbf{\texttt{\_\allowbreak{}\_\allowbreak{}device\_\allowbreak{}\_\allowbreak{} void group\_\allowbreak{}topk\textless{}K\textgreater{}}} --- A templated device function implementing group-level top-K selection using shuffle instructions. It operates on raw pointers (\texttt{f32*}, \texttt{s32*}), not Croqtile spans.

\textbf{\texttt{\_\allowbreak{}\_\allowbreak{}device\_\allowbreak{}\_\allowbreak{} T SHFL\_\allowbreak{}XOR}} --- A lane-shuffle intrinsic wrapper. The \texttt{\_\allowbreak{}\_\allowbreak{}forceinline\_\allowbreak{}\_\allowbreak{}} and \texttt{\_\allowbreak{}\_\allowbreak{}shfl\_\allowbreak{}xor\_\allowbreak{}sync} are device code --- Croqtile passes them through unchanged.

The \co@call warp_topk<8>(...)@ statement supplies the helper's compile-time configuration and ordinary arguments. Croqtile converts spanned-data arguments to element pointers, and the target compiler validates the resulting C++ call.

\textbf{\texttt{inthreads (tid == 0)}} --- Only lane 0 of each group writes back results. Note this uses \texttt{inthreads} without \texttt{.async} --- a sequential thread filter, not a concurrent region.

This pattern --- Croqtile handles the parallelism, tiling, and memory orchestration; \texttt{\_\_device\_\_} functions handle the per-group algorithm --- is common in production kernels like MoE (mixture-of-experts) top-K, custom reductions, and specialized math.
